\documentclass[11pt]{article}
\usepackage[T1]{fontenc}\usepackage[utf8]{inputenc}
\usepackage{lmodern,amsmath,amssymb,amsthm,mathtools,microtype,xcolor}
\usepackage[margin=1in]{geometry}
\usepackage[colorlinks=true,linkcolor=blue!50!black,urlcolor=blue!50!black,citecolor=blue!50!black]{hyperref}
\newtheorem{theorem}{Theorem}[section]\newtheorem{proposition}[theorem]{Proposition}
\newtheorem{lemma}[theorem]{Lemma}\newtheorem{corollary}[theorem]{Corollary}
\theoremstyle{remark}\newtheorem{remark}[theorem]{Remark}
\newcommand{\E}{\mathbb E}\newcommand{\Prb}{\mathbb P}\renewcommand{\Pr}{\mathbb P}\newcommand{\R}{\mathbb R}
\setlength{\emergencystretch}{3em}
\hypersetup{pdftitle={Mean-field Bayesian circuit learning at every width above target complexity},pdfauthor={ChatGPT-6 Astra; Kaarel Hanni; Dmitry Vaintrob}}
\title{Mean-field Bayesian circuit learning\\at every width above target complexity}
\author{ChatGPT-6 Astra \and Kaarel H\"anni \and Dmitry Vaintrob\thanks{Authors are listed in alphabetical order.}}
\date{October 15, 2026}
\begin{document}\maketitle
\begin{abstract}
We prove a sample bound for Bayesian neural-network learning that depends
on the size of a small target implementation and is independent of the
learner's potentially much larger width. The model is a fixed-depth tanh
network at mean-field Gaussian scaling: input-layer weight variance is
$1/d$, internal weight variance is $1/N$
and readout variance is $1/N^2$
at hidden width $N$. Before training, the prior retains only networks whose
outputs have absolute value greater than $1/2$ on the entire Boolean input
cube. Learning then conditions on noiseless labels. The sample bound is
polynomial in input dimension and the target's bounded-weight network
complexity, uniformly at every width at least that complexity. For any
common training and test input distribution, sufficiently many examples
make the posterior, with high probability over the training sample, assign
the correct label probability strictly greater than $2/3$ on all test inputs
outside a set of small probability.
The proof bounds the total prior probability of functions lacking a small
description and compares it with the probability assigned to the target.
Its main step compresses likely hidden-layer covariance patterns into
finite networks. All arguments apply at finite width; no infinite-width
limit is required. The polynomial degree and constants depend on depth,
so the theorem does not provide a bound polynomial in depth.
\end{abstract}

\section{Introduction}
A network with many more parameters than training examples can represent
many functions that fit the data. To explain why Bayesian learning should
nevertheless generalize, one would like a sample bound determined by the
complexity of the target, rather than by the number of parameters available
to the learner. This paper proves such a bound for a Gaussian neural-network
prior at mean-field scaling.

The model is a fully connected tanh network of depth $L$ and hidden width
$N$, with a linear readout. The Gaussian variances are $1/d$ in the input
layer, $1/N$ in internal layers, and $1/N^2$ in the readout. We retain only
networks with absolute output margin greater than $1/2$ at every Boolean
input, then condition on exact agreement with noiseless training labels.
This defines a probability distribution over Boolean functions and, after
training, a posterior distribution over their predictions. The global
margin restriction is part of the prior, including at unobserved inputs.

Suppose the target has a depth-$L$, width-$C$ implementation with that
strict margin and all weights bounded by $C$. Our result gives a sample
bound polynomial in $d,C$, the inverse desired error, and the logarithm of
the inverse failure probability, uniformly for every learner width
$N\geq C$. It holds for any common training and test distribution on the
Boolean cube. Outside a set of prescribed small test probability, the
posterior gives the correct label probability greater than $2/3$.
Section~\ref{sec:model} states the precise guarantee in
terms of the smallest such implementation size $C_L(f)$.

The result is finite-dimensional throughout. It does not require the learner
to be much wider than the target, and it does not take a width limit before
training. The complexity measure is a bounded-weight network complexity at
the specified depth, so no translation from arbitrary Boolean circuits is
implicit. Depth remains fixed in the sample bound: its polynomial degree
and constants may grow rapidly with $L$. As with an Occam learning bound,
the conclusion is statistical and does not establish an efficient method
for computing the posterior.

\paragraph{How excess width is controlled.}
We compare the prior mass of the target with the total mass of functions
that cannot be described by small networks. A small target can be embedded
at any larger width by replicating its neurons; a neighborhood of this
embedding supplies a lower bound on its prior mass. For the competing
functions, we show that all but a sufficiently small amount of prior mass
can be compressed into one finite family whose size is independent of $N$.
Training labels rule out the inaccurate members of this family by an
elementary union bound. This last step explains why a bound on the
\emph{total} exceptional mass is needed, rather than a separate small
probability estimate for each complicated function.

The main difficulty is to construct that finite family at the mean-field
scaling, where the random internal matrices retain operator norm of
constant order as width grows. We work instead with the correlations of
hidden activations across inputs. Finite exponential tests, a convex
duality argument, and covariance stability permit these correlations to be
represented by smaller networks. The proof supplies all three
ingredients and then completes the posterior estimate.

\section{Model and main result}\label{sec:model}
Let $d\geq1$, $L\geq2$, and $\mathcal X=\{-1,1\}^d$.
The augmented input is $\bar x=(1,x_1,\ldots,x_d)$. A depth-$L$ network
has $L-1$ hidden layers of common width $N$ and computes
\[
 h^0(x)=\bar x,\qquad h^s(x)=\tanh(W^sh^{s-1}(x))\quad(1\leq s<L),
 \qquad z_w(x)=W^Lh^{L-1}(x).
\]
There are no additional biases. The product Gaussian prior $\mu_N$ has
independent centered entries with standard deviations
\[
 \sigma_1=d^{-1/2},\qquad
 \sigma_s=N^{-1/2}\quad(2\leq s<L),\qquad
 \sigma_L=N^{-1}.
\]
These are the mean-field conventions used here. In the broader family with
variances $N^{1-\gamma}/d$, $N^{-\gamma}$, and $N^{-\gamma-1}$, they
correspond to $\gamma=1$; no other scaling regime is needed in this paper.
For a Boolean function $g$ write
\[
 A_{N,g}=\{w:g(x)z_w(x)>1/2\text{ for every }x\in\mathcal X\},
 \qquad A_N=\bigcup_g A_{N,g},
 \qquad \pi_N(g)=\frac{\mu_N(A_{N,g})}{\mu_N(A_N)}.
\]
Thus the global margin restriction is imposed before the function prior is
conditioned on training labels. Let $C_L(f)$ be the least integer $C\geq1$
for which a depth-$L$, width-$C$ network implements $f$ with this strict
margin and every weight bounded in absolute value by $C$.
We assume $C_L(f)<\infty$.

For independent inputs $x_1,\ldots,x_n\sim D$, the training sample is
$S_n=((x_i,f(x_i)))_{i=1}^n$, and the posterior is
\[
 \Pi_N(g\mid S_n)=
 \frac{\pi_N(g)\,\mathbf1\{g(x_i)=f(x_i)\ \forall i\}}
 {\sum_h\pi_N(h)\,\mathbf1\{h(x_i)=f(x_i)\ \forall i\}}.
\]
Its average error is
$R_D(\Pi_N)=\E_{g\sim\Pi_N}\Prb_{x\sim D}\{g(x)\ne f(x)\}$.

\begin{theorem}[Learning at every width above target complexity]
\label{thm:main}
For every fixed depth $L\geq2$ there is a polynomial $p_L$ such that the
following holds. Let $C=C_L(f)<\infty$, let $N\geq C$, and let
$\varepsilon,\delta\in(0,1)$. For every distribution $D$ on $\mathcal X$,
if
\begin{equation}\label{eq:main-sample}
 n\geq\frac6\varepsilon\left[
 p_L\!\left(d,C,\log\frac1\varepsilon\right)
       +\log\frac6{\varepsilon\delta}\right],
\end{equation}
then, with probability at least $1-\delta$ over $S_n$,
\[
 R_D(\Pi_N)\leq\varepsilon/3,
 \qquad
 \Prb_{x\sim D}\!\left\{
   \Pi_N(g(x)=f(x)\mid S_n)>2/3\right\}\geq1-\varepsilon.
\]
The polynomial is independent of width. Its degree and constants may depend
on depth. The proof gives explicit quantities that can replace $p_L$.
\end{theorem}

The statement is information-theoretic: efficient posterior computation or
efficient construction of the compressed network is not asserted.
A fixed confidence threshold $\theta\in(1/2,1)$ can replace $2/3$, with the
corresponding change in constants. Thresholds below $2/3$ follow directly.
For $\theta\geq2/3$, apply the theorem with $\varepsilon$ replaced by
$3(1-\theta)\varepsilon$; Markov's inequality then bounds the probability of
$\{\Pi_N(g(x)=f(x)\mid S_n)\leq\theta\}$ by $R_D(\Pi_N)/(1-\theta)\leq\varepsilon$.

\paragraph{Proof roadmap.}
Section~\ref{sec:cloning} constructs a Gaussian neighborhood of the embedded
target with probability at least $e^{-N\operatorname{poly}_L(d,C)}$.
Section~\ref{sec:sparse-dual} develops finite tests for empirical hidden-layer
covariances. These tests identify a family of controlled changes of measure,
called quadratic tilts, whose covariances approximate the observed ones;
the exceptional probability is exponentially small in $N$.
Section~\ref{sec:compression} proves stability of the tilted covariances and
represents them by finite networks. We then choose the approximation budgets
backward through the layers, obtaining a description bound independent of
learner width. The final section compares the exceptional mass with the
target mass and applies a finite-class posterior estimate. A direct
target-mass estimate covers every smaller width where the compression
probability bound is unavailable.

% Standalone section: macros deliberately limited to standard AMS commands.
\section{A target-mass bound at every admissible width}\label{sec:cloning}
The following construction bounds the target's prior cost at every width $N\geq C$,
at the mean-field exponent.

\begin{proposition}[Uniform cloning mass bound]\label{prop:cloning}
Fix $L\geq2$ and an activation $\phi$ satisfying $|\phi|\leq1$ and having Lipschitz constant at most one. There is a function $Q_L(d,C)$ bounded by a polynomial such that for every Boolean target with
$C_L(f)\leq C$, and every $N\geq C$,
\[
 -\frac{2}{N}\log\Pr\{f(x)z_w(x)>1/2\text{ for all }x\}
 \leq Q_L(d,C).
\]
The prior variances and absence of later-layer biases are exactly those in the
main model.
\end{proposition}
\begin{proof}
Choose a width-$C$ teacher, increase its readout by a factor of three, and write
its hidden matrices as $B^s$ and readout as $v$. Thus the teacher has margin
strictly greater than $3/2$, hidden entries bounded by $C$, and readout entries
bounded by $3C$. If the original witness has smaller width, first use its actual
width in the construction below; all estimates only improve.
Partition $\{1,\ldots,N\}$ into $C$ groups of sizes $k_a$ differing by at most
one. Then
\[
 1\leq k_a,\qquad \frac{\max k_a}{\min k_a}\leq2,
 \qquad \frac N{k_a}\leq2C.
\]
Let $u_a=k_a^{-1/2}{\bf1}_{\text{group }a}$ and let $P$ be the orthogonal
projection onto their span. The exact clone $w_*$ has first-layer row $B^1_b$
for a neuron in group $b$, hidden entry $B^s_{ba}/k_a$ from group $a$ to group
$b$, and readout $v_b/k_b$ on group $b$. Every hidden activation of the exact
clone is group-constant. Moreover,
\[
 \|W_*^s\|_{\mathrm{op}}\leq\sqrt2 C^2\quad(2\leq s<L),
 \qquad \sqrt N\|W_*^L\|_2\leq3\sqrt2 C^2.
\]
The first bound follows by restricting to the group-constant subspace: the
resulting matrix has entries $\sqrt{k_b/k_a}B^s_{ba}$, and its Frobenius norm
is at most $\sqrt2 C^2$.

Set $t=1$. The random coordinates
$\sqrt d W^1_{ij}$, $\sqrt N W^s_{ij}$ for $2\leq s<L$, and
$\alpha_j=NW^L_j$ are independent centered Gaussians of standard deviation
$t$. Impose the following cylinder event $\mathcal C_\rho$, for
$0<\rho\leq1$:
\begin{itemize}
\item every $\sqrt d W^1_{ij}$ is within $\rho t$ of
$\sqrt d (W_*^1)_{ij}$;
\item for every hidden row $j$ and group $a$, the coordinate
$\langle\sqrt N W^s_{j,\cdot},u_a\rangle$ is within $\rho t$ of its
cloned value $\sqrt{N/k_a}B^s_{ba}$;
\item every $\alpha_j$ is within $\rho t$ of its cloned value $Nv_b/k_b$.
\end{itemize}
All these constraints concern independent scalar Gaussian coordinates.
The orthogonal coordinates of each internal row remain independent prior
Gaussians. Consequently write
\[
 W^s=W_*^s+R^s+E^s,\qquad R^s=R^sP,\quad E^sP=0.
\]
On the cylinder, $\|R^s\|_{\mathrm{op}}\leq\|R^s\|_F\leq\rho t\sqrt C$.
The free matrix is distributed as $(t/\sqrt N)G^s(I-P)$, with $G^s$ a standard
Gaussian matrix, independently of the cylinder coordinates. Let
\[
 b_L=\sqrt{16\log9+8\log(4L)}.
\]
For completeness, a $1/4$-net of the unit sphere in $\R^N$ has at
most $9^N$ points, by the usual disjoint-ball volume argument. For a matrix
$G$, approximating both unit vectors in its bilinear form by this net gives
$\|G\|_{\mathrm{op}}\leq2\max_{u,v\text{ in the net}}|u^{\mathsf T}Gv|$.
For standard Gaussian $G$, each displayed bilinear form is standard normal.
Its elementary Chernoff bound and a union bound give
\[
 \Prb\{\|G\|_{\mathrm{op}}>b_L\sqrt N\}
 \leq2\cdot9^{2N}e^{-b_L^2N/8}
 =2(4L)^{-N}\leq(2L)^{-1}.
\]
Thus a union bound implies, conditionally on
$\mathcal C_\rho$, that all $\|E^s\|_{\mathrm{op}}\leq b_Lt$ with probability
at least $1/2$. This is also valid when $L=2$, where there are no such matrices.

Let $e_s(x)=N^{-1/2}\|h^s_w(x)-h^s_*(x)\|_2$. The input constraints give
$e_1(x)\leq\rho(d+1)/\sqrt d\leq\rho(d+1)$. Since $E^sh^{s-1}_*=0$,
boundedness and Lipschitz continuity give, simultaneously for every input,
\[
 e_s(x)\leq (\sqrt2 C^2+\sqrt C+b_L)e_{s-1}(x)+\rho\sqrt C.
\]
Put $A=2+3C^2+b_L$. Induction yields
$e_s(x)\leq\rho(d+L+1)A^{L-2}$ for $s<L$. At the output,
\[
 |z_w(x)-z_*(x)|
 \leq3\sqrt2 C^2 e_{L-1}(x)+\rho.
\]
Hence choosing
\[
 \rho=\bigl[2+9C^2(d+L+1)A^{L-2}\bigr]^{-1}
\]
ensures output discrepancy at most $1/2$. All networks in the cylinder with
the operator-norm event therefore implement $f$ with margin greater than one.

It remains to price the cylinder. Its number of constrained coordinates is
\[
 m=N\bigl[d+2+(L-2)C\bigr].
\]
The sum of their squared centers is at most $NB_0$, where one may take
\[
 B_0=d(d+1)C^2+2(L-2)C^4+36C^4.
\]
Indeed, a first-layer center has magnitude at most $\sqrt d C$, an internal
center at most $\sqrt{2C}C$, and a readout center at most $6C^2$.
For $Z\sim N(0,t^2)$ and $|Z-a|\leq\rho t$,
\[
 \Pr\{|Z-a|\leq\rho t\}
 \geq\frac{2\rho}{\sqrt{2\pi}}
       \exp\{-a^2/t^2-\rho^2\}.
\]
Multiplication and $t=1$ show
\[
 -\log\Pr(\mathcal C_\rho)
 \leq N B_0
 +N[d+2+(L-2)C]\log\frac{e\sqrt{2\pi}}{2\rho}.
\]
The desired probability is at least $\Pr(\mathcal C_\rho)/2$. Since $N\geq1$, this proves the claim with
\[
 Q_L(d,C)=2B_0+\frac{4[d+2+(L-2)C]}{\rho}+2.
\]
Here we used $\log(e\sqrt{2\pi}/(2\rho))\leq2/\rho$ for $0<\rho\leq1$.
The displayed expression is a polynomial in $d,C$ at every fixed depth.
\end{proof}

% Self-contained proof fragment. Requires amsmath, amssymb, amsthm and
% theorem environments lemma/proposition. No references to older manuscripts.
\section{A finite family of entropy certificates}\label{sec:sparse-dual}
The following bound replaces a discretization of the full feature space by
sparse linear tests. Its logarithmic dependence on the number of observables
is essential when those observables are indexed by all pairs of Boolean
inputs.

Let $(\Omega,\mathcal A,\mu)$ be any probability space, let
$F=(F_1,\ldots,F_D):\Omega\to[-1,1]^D$ be measurable, and write
\[
 \Lambda(\lambda)=\log\int e^{\lambda\cdot F}\,d\mu,
 \qquad
 \Psi_G(q)=\max_{\|\lambda\|_1\leq G}
       \{\lambda\cdot q-\Lambda(\lambda)\},\qquad G>0.
\]
The maximum exists: its domain is compact and its objective is continuous.
The norm $\|\lambda\|_1$ here is the sum of the absolute values of its
$D$ coordinates.

\begin{lemma}[A bounded exponential tilt with controlled error]
\label{lem:bounded-dual-lift}
For every $q\in\mathbb R^D$, there is a maximizer $\lambda_*$ in the
preceding display such that the probability measure
\[
 \frac{d\nu_*}{d\mu}
     =\exp\{\lambda_*\cdot F-\Lambda(\lambda_*)\}
\]
satisfies the exact identity
\begin{equation}\label{eq:dual-lift-identity}
 \Psi_G(q)
   =D(\nu_*\Vert\mu)
      +G\left\|q-\int F\,d\nu_*\right\|_\infty.
\end{equation}
Consequently, if $\Psi_G(q)\leq K$, then
\[
 D(\nu_*\Vert\mu)\leq K,
 \qquad
 \left\|q-\int F\,d\nu_*\right\|_\infty\leq K/G.
\]
\end{lemma}
\begin{proof}
Boundedness of $F$ permits differentiation under the integral and gives
$\nabla\Lambda(\lambda_*)=\int F\,d\nu_*=:m_*$. The first-order
condition for a differentiable concave function maximized over the convex
$\ell^1$ ball is
\[
 (q-m_*)\cdot(\lambda-\lambda_*)\leq0
       \qquad(\|\lambda\|_1\leq G).
\]
Taking the maximum of the left linear functional over that ball gives
\[
 \lambda_*\cdot(q-m_*)=G\|q-m_*\|_\infty.
\]
Moreover,
$D(\nu_*\Vert\mu)=\lambda_*\cdot m_*-\Lambda(\lambda_*)$.
Adding these two identities proves \eqref{eq:dual-lift-identity}.
Relative entropy is nonnegative, which proves the stated consequences.
The log density of $\nu_*$ relative to $\mu$ is bounded, so its entropy
is finite. No support, density, or nonsingularity assumption on $\mu$
was used.
\end{proof}

For an integer $s\geq1$, introduce the finite set
\[
 \mathcal V_{G,s}
 =\left\{\frac1s\sum_{j=1}^s a_j:
        a_j\in\{0,Ge_1,-Ge_1,\ldots,Ge_D,-Ge_D\}\right\}.
\]
Thus $\#\mathcal V_{G,s}\leq(2D+1)^s$ and every vector in this set has
$\ell^1$ norm at most $G$.

\begin{lemma}[Sparse approximation of the dual value]
\label{lem:sparse-dual-net}
For every $q\in\mathbb R^D$,
\begin{equation}\label{eq:sparse-dual-net}
 \Psi_G(q)\leq\frac{G^2}{2s}
   +\max_{v\in\mathcal V_{G,s}}
        \{v\cdot q-\Lambda(v)\}.
\end{equation}
\end{lemma}
\begin{proof}
Fix $\lambda$ with $\|\lambda\|_1\leq G$. Let the random vector $A$
take the value $G\operatorname{sign}(\lambda_i)e_i$ with probability
$|\lambda_i|/G$ for each nonzero coordinate, and take value zero with the
remaining probability. Then $\mathbb EA=\lambda$.
Take independent copies $A_1,\ldots,A_s$ and put
$V=s^{-1}\sum_j A_j\in\mathcal V_{G,s}$.
For each fixed $\omega\in\Omega$, the scalar $A\cdot F(\omega)$ lies
in $[-G,G]$. Hoeffding's exponential-moment inequality therefore gives
\[
 \mathbb E\exp\{(V-\lambda)\cdot F(\omega)\}
 \leq\exp\{G^2/(2s)\}.
\]
Tonelli's theorem and Jensen's inequality imply
\begin{align*}
 \mathbb E\Lambda(V)
 &=\mathbb E\log\int e^{V\cdot F}\,d\mu\\
 &\leq\log\mathbb E\int e^{V\cdot F}\,d\mu\\
 &\leq\Lambda(\lambda)+G^2/(2s).
\end{align*}
Since $\mathbb E(V\cdot q)=\lambda\cdot q$, at least one realization
$v\in\mathcal V_{G,s}$ satisfies
\[
 v\cdot q-\Lambda(v)
 \geq\lambda\cdot q-\Lambda(\lambda)-G^2/(2s).
\]
Take $\lambda$ to be a maximizer defining $\Psi_G(q)$.
\end{proof}

\begin{proposition}[A nonasymptotic entropy certificate]
\label{prop:finite-entropy-certificate}
Let $X_1,\ldots,X_N$ be independent with law $\mu$, and put
$q_N=N^{-1}\sum_{i=1}^N F(X_i)$. For every $K>0$, $G>0$, and integer
$s\geq1$,
\begin{equation}\label{eq:finite-entropy-tail}
 \Pr\{\Psi_G(q_N)>K\}
 \leq\exp\left\{s\log(2D+1)
          -N\left(K-\frac{G^2}{2s}\right)\right\}.
\end{equation}
Outside this exceptional event, the empirical moment vector is within
$K/G$ in maximum norm of the moments of an exponential tilt with parameter
$\ell^1$ norm at most $G$ and relative entropy at most $K$.
In particular, choosing $s=\lceil G^2\rceil$ bounds the right side by
$\exp\{\lceil G^2\rceil\log(2D+1)-N(K-1/2)\}$.
\end{proposition}
\begin{proof}
For each fixed $v\in\mathcal V_{G,s}$, independence gives exactly
\[
 \mathbb E\exp\{N[v\cdot q_N-\Lambda(v)]\}=1.
\]
Markov's inequality bounds the probability that its bracket exceeds
$K-G^2/(2s)$ by $\exp\{-N(K-G^2/(2s))\}$.
The sparse approximation lemma and a union bound over
$\mathcal V_{G,s}$ prove \eqref{eq:finite-entropy-tail}.
The bounded-dual lift lemma gives the final assertion.
\end{proof}

\paragraph{Application to hidden-layer covariances.}
Take $m$ input points, $D=m^2$, $\mu=\mathcal N(0,Q)$ for any positive
semidefinite covariance $Q$, and $F_{ij}(z)=\tanh(z_i)\tanh(z_j)$,
indexed by all ordered pairs $(i,j)$. Reshape a maximizing parameter
$\lambda_*$ into a matrix and replace it by its symmetric part $\Gamma$.
This preserves its quadratic form and does not increase its entrywise
$\ell^1$ norm. The lifted covariance is then exactly
\[
 \mathsf F_\Gamma(Q)=
 \frac{\mathbb E_{Z\sim\mathcal N(0,Q)}
    [\tanh(Z)\tanh(Z)^\top e^{\tanh(Z)^\top\Gamma\tanh(Z)}]}
      {\mathbb E_{Z\sim\mathcal N(0,Q)}
    e^{\tanh(Z)^\top\Gamma\tanh(Z)}}.
\]
Thus, except with the probability in \eqref{eq:finite-entropy-tail}, the
empirical activation covariance $\widehat Q$ satisfies
\[
 \|\widehat Q-\mathsf F_\Gamma(Q)\|_{\max}\leq K/G,
 \qquad \|\Gamma\|_{1,\mathrm{ent}}\leq G.
\]
This assertion is valid conditionally on any preceding network layers,
using their empirical covariance as $Q$. Its probability bound is uniform
in that covariance, including singular covariances. On the Boolean cube,
$m=2^d$ and the logarithmic factor is
$\log(2m^2+1)\leq(2d+2)\log2$, which is linear in $d$.

\section{Covariance stability and finite realization}\label{sec:compression}
We first bound how a quadratic tilt changes when its Gaussian covariance
changes, then realize repeated tilted covariance maps by a finite network.
Write $\phi=\tanh$. Recall that $\|A\|_{\max}=\max_{ij}|A_{ij}|$ and
$\|A\|_{1,\mathrm{ent}}=\sum_{ij}|A_{ij}|$ for matrices on the finite input
set. These are different norms; their dual pairing is what avoids a factor
of $2^d$ below.
For a positive semidefinite covariance $Q$ and a symmetric matrix $\Gamma$,
let
\[
 \mathsf F_\Gamma(Q)=
 \frac{\mathbb E_{Z\sim\mathcal N(0,Q)}
 [\phi(Z)\phi(Z)^\top\exp\{\phi(Z)^\top\Gamma\phi(Z)\}]}
 {\mathbb E_{Z\sim\mathcal N(0,Q)}
 [\exp\{\phi(Z)^\top\Gamma\phi(Z)\}]}.
\]
All activations are coordinatewise. The definition is valid for singular $Q$.

\begin{lemma}[A dimension-independent covariance comparison]\label{lem:tilted-covariance}
If $g=\|\Gamma\|_{1,\mathrm{ent}}$, then for all positive semidefinite
$Q,Q'$ of the same size,
\[
 \|\mathsf F_\Gamma(Q)-\mathsf F_\Gamma(Q')\|_{\max}
 \leq 8(1+g)^2\|Q-Q'\|_{\max}.
\]
\end{lemma}
\begin{proof}
Put $V(z)=\phi(z)^\top\Gamma\phi(z)$ and
$Q_t=(1-t)Q+tQ'$. For a twice differentiable bounded test function $F$, let
$\mathbb E_t$ denote expectation under the Gaussian of covariance $Q_t$
reweighted by $e^V$. Gaussian differentiation and the quotient rule give
\begin{align*}
 \frac{d}{dt}\mathbb E_tF
 =\frac12\sum_{a,b}(Q'-Q)_{ab}\bigl\{
 &\mathbb E_t[\partial_{ab}F+
       \partial_aF\partial_bV+\partial_bF\partial_aV]\\
 &+\operatorname{Cov}_t(F,\partial_{ab}V+
                                   \partial_aV\partial_bV)\bigr\}.
\end{align*}
This formula has no reciprocal partition-function bound: all terms are
expectations under a probability measure. Define
$M_1(H)=\sup_z\sum_a|\partial_aH(z)|$ and
$M_2(H)=\sup_z\sum_{a,b}|\partial_{ab}H(z)|$.
Because $|\phi|,|\phi'|\leq1$ and $|\phi''|\leq2$,
\[
 M_1(V)\leq2g,\qquad M_2(V)\leq6g.
\]
For $F(z)=\phi(z_i)\phi(z_j)$, including $i=j$,
$\|F\|_\infty\leq1$, $M_1(F)\leq2$, and $M_2(F)\leq6$.
The absolute derivative is therefore at most
\[
 \frac12\|Q-Q'\|_{\max}
 \{6+8g+2(6g+4g^2)\}
 =(3+10g+4g^2)\|Q-Q'\|_{\max}.
\]
Integrating in $t$ proves the assertion, with the displayed slightly larger
constant. One may initially replace both covariances by $Q+\epsilon I$ and
$Q'+\epsilon I$. The functions and all derivatives used above are bounded,
and the partition functions lie between $e^{-g}$ and $e^g$, so weak Gaussian
convergence lets $\epsilon\downarrow0$. This also proves the singular case.
\end{proof}

Applying one tilted covariance map per hidden layer produces a sequence
of matrices, which we call a \emph{tilted tower}. The next proposition
shows that the output specified by such a sequence has a small network
implementation. The tilt parameters select a network whose existence is
proved by sampling; they are not additional inputs to that network.

\begin{proposition}[Finite realization of a bounded-tilt tower]\label{prop:tilt-compress}
Fix $L\geq2$, $K\geq1$, and $G\geq1$. Let $Q_0$ be the covariance of the
Gaussian first-layer field $w\cdot\bar x$, with
$w\sim\mathcal N(0,I_{d+1}/d)$, indexed by $x\in\mathcal X$.
Suppose symmetric matrices $\Gamma_1,\ldots,\Gamma_{L-1}$ and a vector
$\beta\in\mathbb R^{\mathcal X}$ satisfy
\[
 \|\Gamma_s\|_{1,\mathrm{ent}}\leq G,\qquad
 \|\beta\|_1\leq4K,
 \qquad Q_s=\mathsf F_{\Gamma_s}(Q_{s-1}).
\]
Then the function $z(x)=(Q_{L-1}\beta)_x$ can be uniformly approximated to
accuracy $1/8$ by a depth-$L$ network whose width and weights are bounded by
$\operatorname{poly}_L(d,K,G)$.
If $f(x)z(x)\geq7/8$ everywhere, the finite network implements $f$ with margin at least $3/4$.
\end{proposition}
\begin{proof}
Let $b=8(1+G)^2$ and
\[
 \eta=[32K(L-1)b^{L-2}]^{-1}.
\]
Choose the common width $n$ to be the least integer satisfying
\[
 n\geq2\eta^{-2}\log\bigl(8(L-1)\,2^{2d}\bigr).
\]
The construction is random only to prove existence.
The matrices $\Gamma_s$ may be indexed by the entire cube; they select the
sampling laws, but are not stored or evaluated by the resulting finite
network. No efficient sampling or construction algorithm is claimed.
At the first layer, independently sample $n$ Gaussian weight vectors
$w_i\sim\mathcal N(0,I_{d+1}/d)$ tilted by
\[
 \exp\{h(w_i)^\top\Gamma_1h(w_i)\},\qquad
 h(w_i)_x=\phi(w_i\cdot\bar x).
\]
Their activation vectors have exactly the law defining
$Q_1=\mathsf F_{\Gamma_1}(Q_0)$.
At each subsequent hidden layer, conditional on the already sampled previous
activation vectors $h_1,\ldots,h_n$, independently sample $n$ standard
Gaussian row vectors $\xi_j\in\mathbb R^n$ tilted by
\[
 \exp\{\phi(Z_j)^\top\Gamma_s\phi(Z_j)\},
 \qquad Z_j(x)=\frac1{\sqrt n}\sum_{i=1}^n\xi_{ji}h_i(x).
\]
Use the actual network weights $\xi_{ji}/\sqrt n$.
If $\widehat Q_{s-1}=n^{-1}\sum_i h_ih_i^\top$ is the preceding empirical
activation covariance, then, conditionally, $Z_j$ before tilting is Gaussian
with exactly that covariance. Thus the new activation vectors are
conditionally independent, with second moment
$\mathsf F_{\Gamma_s}(\widehat Q_{s-1})$.

Every covariance entry is the empirical mean of $n$ variables in $[-1,1]$.
Hoeffding's inequality, a union bound over $2^{2d}$ entries, and conditioning
at successive layers show, with probability at least $3/4$, that
\[
 \|\widehat Q_s-
          \mathsf F_{\Gamma_s}(\widehat Q_{s-1})\|_{\max}\leq\eta
 \quad(1\leq s<L),
\]
where $\widehat Q_0=Q_0$.
Lemma~\ref{lem:tilted-covariance} then implies
\[
 \|\widehat Q_{L-1}-Q_{L-1}\|_{\max}
 \leq (L-1)b^{L-2}\eta.
\]
For each last-layer activation vector $h_i$, choose the deterministic
readout weight $v_i=(\beta\cdot h_i)/n$. The resulting output is
$\widehat z(x)=(\widehat Q_{L-1}\beta)_x$, and therefore
\[
 \|\widehat z-z\|_\infty
 \leq4K(L-1)b^{L-2}\eta=1/8.
\]
Although $\beta$ is indexed by the entire input cube, its only role in the
finite network is to select these $n$ real readout weights; each satisfies
$|v_i|\leq4K/n$. It is not an additional computational input to the network.

It remains to bound the sampled weights without losing polynomial size.
For every sampled row, its tilted density relative to its Gaussian reference
is at most $e^{2G}$, since the quadratic tilt lies in $[-G,G]$.
Let $D_n=n(d+1)+(L-2)n^2$ be the total number of incoming scalar coordinates,
and choose
\[
 B=\bigl[4G+2\log(8\max\{D_n,1\})\bigr]^{1/2}.
\]
Each raw incoming coordinate has reference variance at most one, so, also
conditionally on all preceding layers, its probability of exceeding $B$ in
absolute value is at most $2e^{2G-B^2/2}$.
A union bound gives probability at most $1/4$ that any incoming coordinate
exceeds $B$. The actual internal weights divide the raw coordinates by
$\sqrt n$ and so satisfy the same bound. Intersecting this event with the
covariance-approximation event leaves probability at least $1/2$.
Fix a realization in the intersection. Both $n$ and
$\max\{B,4K/n\}$ are bounded by
polynomials in $d,K,G$ at fixed depth. The stated margin follows from $7/8-1/8=3/4$.
\end{proof}

\begin{lemma}[Counting finite realizations]\label{lem:tilt-count}
Fix $L$, width $n$, and an entrywise bound $B\geq1$. The Boolean functions
having depth-$L$, width-$n$ implementations with margin at least $3/4$ and
weights bounded by $B$ number at most
\[
 \exp\{P\log(3+16B\Lambda)\},\qquad
 P=n(d+1)+(L-2)n^2+n,
 \quad \Lambda=(d+Ln)[1+n(B+1)]^{L-1}.
\]
\end{lemma}
\begin{proof}
Round every weight to a multiple of $(8\Lambda)^{-1}$.
The first hidden layer's coordinatewise error is at most $(d+1)/(8\Lambda)$.
Every subsequent hidden layer has error at most $n(B+1)$ times the preceding
error plus $n/(8\Lambda)$; the same estimate applies at the linear readout.
Induction bounds the final error by $1/8$. Each rounded coordinate has at
most $3+16B\Lambda$ possible values, so the $P$-parameter grid has the stated
size. Rounding preserves every represented Boolean sign function.
\end{proof}


\section{Likely full-margin draws have small descriptions}
We now combine the covariance tests and finite realizations. The goal is
one finite family of Boolean functions that contains the sign function of
every globally margin-separated draw outside a controlled exceptional
event. Its size will be bounded by rounding the weights of the smaller
networks, independently of the learner's width.
Let $m=2^d$ and $r=L-1$. Index fields and covariance matrices by the
$m$ cube inputs. Define
\[
 Q_0(x,y)=\bar x\cdot\bar y/d,
 \qquad
 Q_{s,N}=\frac1N\sum_{i=1}^N h_i^s(h_i^s)^{\mathsf T}
 \quad(1\leq s\leq r),
\]
where $h_i^s$ denotes the activation of neuron $i$ evaluated on all inputs.
The activation matrices in this display are uncentered second moments.
For each internal hidden layer, the preceding second-moment matrix is
exactly the covariance of its centered Gaussian preactivation vector.
Write $\mathsf M_Q$ for the law of $\tanh(Z)$ when
$Z\sim\mathcal N(0,Q)$. Conditional on the preceding layers, the rows
$h_1^s,\ldots,h_N^s$ are independent with law
$\mathsf M_{Q_{s-1,N}}$, with $Q_{0,N}=Q_0$.
This is exact: an internal Gaussian row has preactivation covariance
$N^{-1}\sum_i h_i^{s-1}(h_i^{s-1})^{\mathsf T}$.
No covariance matrix is assumed nonsingular.

For a covariance $\widehat Q$ and a probability law $\mu$ on
$[-1,1]^m$, define
\[
 \mathcal J_G(\widehat Q;\mu)=
 \max_{\|\Gamma\|_{1,\mathrm{ent}}\leq G}
 \left\{\langle\Gamma,\widehat Q\rangle
       -\log\int e^{h^{\mathsf T}\Gamma h}\,d\mu(h)\right\},
 \quad \langle A,B\rangle=\sum_{i,j}A_{ij}B_{ij}.
\]
This is the bounded-observable dual from Section~\ref{sec:sparse-dual}
with $D=m^2$. Symmetrizing a maximizing $\Gamma$ leaves its quadratic
form unchanged and cannot increase its entrywise norm.

Fix $K\geq4$. Choose the layer budgets backward:
\begin{equation}\label{eq:backwards-budgets}
 G_r=32rK^2,\qquad
 G_s=8G_{s+1}(1+G_{s+1})^2\quad(s=r-1,\ldots,1),
 \qquad b_s=8(1+G_s)^2.
\end{equation}
Then $G_s=32rK^2\prod_{j=s+1}^r b_j$ and
$G_1\geq\cdots\geq G_r\geq1$.
Set
\begin{equation}\label{eq:polynomial-threshold}
 q_s=\lceil G_s^2/2\rceil,
 \qquad \ell=\log(2m^2+1),
 \qquad T_K=\left\lceil\max\left\{1,\frac{2q_1\ell}{K-1}\right\}\right\rceil.
\end{equation}
At fixed depth, $G_1$, $q_1$, and $T_K$ are polynomially bounded in $d,K$;
indeed $\ell\leq(2d+2)\log2$ and the cubic recursion gives
\[
 G_1=O_L(K^{2\cdot3^{L-2}}),\qquad
 T_K=O_L(dK^{4\cdot3^{L-2}-1}).
\]

\begin{lemma}[A likely set of parameter draws]\label{lem:good-event}
Let $\mathcal E_K$ be the event that
\[
 \mathcal J_{G_s}(Q_{s,N};\mathsf M_{Q_{s-1,N}})\leq K
 \quad\text{for all }s,
 \qquad N\|W^L\|_2^2\leq K.
\]
For every $N\geq T_K$,
\begin{equation}\label{eq:good-event-tail}
 \mu_N(\mathcal E_K^c)\leq L e^{-NK/8}.
\end{equation}
\end{lemma}
\begin{proof}
At hidden layer $s$, apply Proposition~\ref{prop:finite-entropy-certificate}
conditionally on the preceding layers, with its sparsity parameter equal to
$q_s$. Since
$G_s^2/(2q_s)\leq1$ and $q_s\leq q_1$, each hidden-layer failure
probability is at most
\[
 \exp\{q_1\ell-N(K-1)\}
 \leq e^{-N(K-1)/2}\leq e^{-NK/4}.
\]
The bound is uniform in the preceding covariance, so averaging preserves it.
For the readout, $N^2\|W^L\|_2^2$ has a chi-square distribution with
$N$ degrees of freedom. Its exponential moment at $1/4$ gives
\[
 \Prb\{N\|W^L\|_2^2>K\}
 \leq e^{-NK/4}2^{N/2}\leq e^{-NK/8},
\]
where $K\geq4$ is sufficient. A union bound over the $r$ hidden layers
and the readout proves the claim.
\end{proof}

\begin{lemma}[Minimum readout and a bounded dual vector]
\label{lem:readout-dual}
Suppose $H$ is any $m\times N$ feature matrix, $g\in\{-1,1\}^m$, and
$v\in\R^N$ satisfies $g_i(Hv)_i>1/2$ for every $i$ and
$N\|v\|_2^2\leq K$. Put $Q=HH^{\mathsf T}/N$.
There exists $\beta\in\R^m$ with
\[
 \|\beta\|_1\leq2K,
 \qquad g_i(Q\beta)_i\geq1/2\quad\text{for every }i.
\]
This holds even when $Q$ is singular.
\end{lemma}
\begin{proof}
Minimize $N\|u\|_2^2/2$ subject to $g_i(Hu)_i\geq1/2$.
A minimizer exists by coercivity and closedness, and the given $v$ is a
strictly feasible point. The elementary convex quadratic-program
optimality conditions therefore give multipliers $\lambda_i\geq0$ with
\[
 Nu=H^{\mathsf T}\beta,
 \qquad\beta_i=g_i\lambda_i,
 \qquad\lambda_i[g_i(Hu)_i-1/2]=0.
\]
Consequently $Hu=Q\beta$ and
\[
 \frac12\|\beta\|_1
 =\beta^{\mathsf T}Hu
 =N\|u\|_2^2\leq N\|v\|_2^2\leq K.
\]
This proves the result without inverting $Q$.
\end{proof}

\begin{proposition}[Finite descriptions on the likely set]
\label{prop:good-compression}
For each $K\geq4$ there is a finite class $\mathcal H_K$ of Boolean
functions, independent of $N$, such that every full-margin draw in
$\mathcal E_K$ has its sign function in $\mathcal H_K$, and
\[
 \log|\mathcal H_K|\leq H_L(d,K),
\]
where $H_L$ is polynomially bounded at fixed depth.
\end{proposition}
\begin{proof}
Fix a draw in $\mathcal E_K$. Lemma~\ref{lem:bounded-dual-lift} gives
symmetric matrices $\Gamma_s$ with
$\|\Gamma_s\|_{1,\mathrm{ent}}\leq G_s$ and
\[
 \|Q_{s,N}-\mathsf F_{\Gamma_s}(Q_{s-1,N})\|_{\max}\leq K/G_s.
\]
Construct a genuine tilted tower from these same matrices by
$Q_0^*=Q_0$ and $Q_s^*=\mathsf F_{\Gamma_s}(Q_{s-1}^*)$.
The matrices may depend on the entire draw; only their uniform norm bounds
will be used to count the resulting functions.
For $e_s=\|Q_{s,N}-Q_s^*\|_{\max}$,
Lemma~\ref{lem:tilted-covariance} gives
\[
 e_0=0,\qquad e_s\leq K/G_s+b_s e_{s-1}.
\]
The backward choice in \eqref{eq:backwards-budgets} yields
\begin{equation}\label{eq:tower-approximation}
 e_r\leq\sum_{s=1}^r\frac K{G_s}\prod_{j=s+1}^r b_j
      =\frac1{32K}.
\end{equation}
Using a common budget at every layer would not give this estimate: the
larger tilt also enlarges the next layer's covariance Lipschitz constant.

If the draw has full-margin sign $g$, apply Lemma~\ref{lem:readout-dual}
to its final hidden feature matrix and original readout. This gives
$\|\beta\|_1\leq2K$ and $g_i(Q_{r,N}\beta)_i\geq1/2$.
By \eqref{eq:tower-approximation},
\[
 \|Q_r^*\beta-Q_{r,N}\beta\|_\infty\leq1/16.
\]
Thus $2Q_r^*\beta$ has signed margin at least $7/8$.
Its dual vector $2\beta$ has $\ell^1$ norm at most $4K$, while all its
tilt matrices have entrywise norm at most $G_1$.
Proposition~\ref{prop:tilt-compress} supplies a finite network within
$1/8$ of this output, and hence with signed margin at least $3/4$.
Its width and weights have uniform bounds polynomial in $d,K,G_1$.
Since $G_1$ is polynomial in $K$ at fixed depth, Lemma~\ref{lem:tilt-count}
provides the required finite class and count.
\end{proof}

Here is one explicit choice of the count. Write $G=G_1$,
$b=8(1+G)^2$, and
\begin{align}
 \eta&=[32Kr b^{r-1}]^{-1},&
 n_*&=\left\lceil2\eta^{-2}\log(8r\,2^{2d})\right\rceil,\notag\\
 D_*&=n_*(d+1)+(L-2)n_*^2,&
 B_*&=\max\{1,\sqrt{4G+2\log(8\max\{D_*,1\})},4K/n_*\},\notag\\
 P_*&=D_*+n_*,&
 \Lambda_*&=(d+Ln_*)[1+n_*(B_*+1)]^{L-1}.
 \label{eq:explicit-compression}
\end{align}
Then
\begin{equation}\label{eq:explicit-count}
 H_L(d,K)=P_*\log(3+16B_*\Lambda_*)
\end{equation}
is valid. It counts all sign functions realized by the common bounded
network family, rather than choices of the input-indexed tilt matrices.

\begin{corollary}[Aggregate function-mass bound]\label{cor:aggregate}
For every $N\geq T_K$,
\[
 \mu_N\!\left(\bigcup_{g\notin\mathcal H_K}A_{N,g}\right)
 \leq L e^{-NK/8}.
\]
\end{corollary}
\begin{proof}
The union in question is contained in $\mathcal E_K^c$ by
Proposition~\ref{prop:good-compression}. Apply Lemma~\ref{lem:good-event}.
\end{proof}

\section{The posterior guarantee at every width}
We finish with elementary posterior estimates, including the widths below
the polynomial threshold $T_K$.

\begin{lemma}[A finite relevant class]\label{lem:finite-posterior}
Let a prior on a finite function class assign positive mass to a target $f$.
Suppose $f\in\mathcal H$, $\log|\mathcal H|\leq H$, and
$\pi(\mathcal H^c)/\pi(f)\leq\tau$.
If $n\geq a^{-1}[H+\log(1/\delta)]$, then, with sample probability
at least $1-\delta$, its consistency posterior has average error at most
$a+\tau$.
\end{lemma}
\begin{proof}
A fixed function of true error greater than $a$ survives $n$ independent
labels with probability at most $(1-a)^n\leq e^{-na}$.
A union bound excludes all such survivors in $\mathcal H$ except on an
event of probability at most $e^{H-na}$.
For every sample the posterior denominator is at least $\pi(f)$, so its
mass outside $\mathcal H$ is at most $\tau$.
\end{proof}

\begin{proof}[Proof of Theorem~\ref{thm:main}]
Let $A=Q_L(d,C)$ from Proposition~\ref{prop:cloning}, so $A\geq1$ and
\[
 \mu_N(A_{N,f})\geq e^{-AN/2}\quad\text{for every }N\geq C.
\]
Set
\begin{equation}\label{eq:learning-budget}
 K=4A+8\log(6L/\varepsilon),
 \qquad
 P_L(d,C,\varepsilon)=
 \max\{H_L(d,K),\,A T_K/2\}.
\end{equation}
These are explicit, and $P_L$ is polynomially bounded in
$d,C,\log(1/\varepsilon)$ at fixed depth.
We show it suffices to replace $p_L$ in \eqref{eq:main-sample} by $P_L$.

First let $N\geq T_K$. Corollary~\ref{cor:aggregate}, the target-mass
bound, and cancellation of the global-margin normalization give
\[
 \frac{\pi_N(\mathcal H_K^c)}{\pi_N(f)}
 \leq L\exp\{-N(K/8-A/2)\}
 =L\exp\{-N\log(6L/\varepsilon)\}\leq\varepsilon/6.
\]
Necessarily $f\in\mathcal H_K$, since otherwise this ratio would be at
least one. Lemma~\ref{lem:finite-posterior}, with
$a=\tau=\varepsilon/6$, gives $R_D(\Pi_N)\leq\varepsilon/3$ with
probability at least $1-\delta$ as soon as
\[
 n\geq\frac6\varepsilon[H_L(d,K)+\log(1/\delta)].
\]

Now let $C\leq N<T_K$. The target-mass bound and $\mu_N(A_N)\leq1$
give $-\log\pi_N(f)\leq AN/2\leq A T_K/2$.
Let $B_n$ be the posterior mass on functions of true error greater than
$\varepsilon/6$. Using the target as a lower bound on the posterior
normalizer yields
\[
 \E_{S_n}B_n
 \leq\frac{1}{\pi_N(f)}
       \sum_{g:\operatorname{err}_D(g)>\varepsilon/6}
       \pi_N(g)(1-\operatorname{err}_D(g))^n
 \leq\exp\{A T_K/2-n\varepsilon/6\}.
\]
Markov's inequality gives $B_n\leq\varepsilon/6$ with probability at
least $1-\delta$ if
\[
 n\geq\frac6\varepsilon
       [A T_K/2+\log(6/(\varepsilon\delta))].
\]
The remaining posterior mass has error at most $\varepsilon/6$, proving
$R_D(\Pi_N)\leq\varepsilon/3$ in this width range as well.
Both sufficient sample bounds follow from \eqref{eq:main-sample} with
$P_L$ in place of $p_L$. A polynomial majorant gives the stated theorem.

Finally, at every input whose correct-label posterior probability is at
most $2/3$, the posterior error probability is at least $1/3$.
Consequently the probability of this failure event is at most
$3R_D(\Pi_N)\leq\varepsilon$. This proves the strict $>2/3$ success
criterion without excluding equality from the failure event.
\end{proof}

\begin{remark}[Activation assumptions]
Tanh is the activation used in the theorem and throughout the model.
The proof also applies to any bounded twice continuously differentiable
activation satisfying $|\phi|\leq1$, $|\phi'|\leq1$, and
$|\phi''|\leq2$, provided the target has the stated finite bounded-weight
implementation. The entropy tests use only bounded features; the covariance
comparison uses precisely these derivative bounds. Analyticity,
invertibility, and oddness are not used. Other finite derivative bounds
change the constants.
\end{remark}

The proof establishes a width-uniform statistical guarantee with explicit,
though large, polynomial bounds at fixed depth. In particular, it does not
assert useful numerical sample sizes or uniform polynomial degrees as depth
grows. The finite-width entropy certificates and the backward layer budgets
are what permit the same theorem to cover every $N\geq C_L(f)$.
\end{document}
